\section{A differentiated explicit formula}
For each character \(\chi\bmod q\) let \(\chi^*\) be its primitive inducer. Write \(m_\chi(\gamma)\) for the multiplicity at \(1/2+i\gamma\) of \(L(s,\chi^*)\), taking zero for absence. GRH places every nontrivial zero on this line. Let \(\Gamma\) be the union of the distinct positive ordinates and define
\[
 b_\gamma(a)=-\frac1d\sum_{\chi\bmod q}\overline{\chi(a)}\,m_\chi(\gamma)
 \qquad(\gamma\in\Gamma\cup(-\Gamma)\cup\{0\}).
\]
Conjugating the characters gives \(b_{-\gamma}=\overline{b_\gamma}\), and \(b_0\) is real. Put \(r_q(a)=|\{b\in G:b^2=a\}|\) and
\[
 F(t)=b_0+2\operatorname{Re}\sum_{\gamma\in\Gamma}b_\gamma\e^{-h\gamma^2}\e^{i\gamma t}.
\]

\begin{proposition}\label{pr:explicit}
For each fixed derivative order \(k\ge0\), as \(t\to\infty\),
\begin{equation}
 (\QQ-F)^{(k)}(t)=O_{q,h,k}(\e^{-t/4}),\qquad
 \left(\PP-F+\frac{r_q}{2d}\right)^{(k)}(t)=o_{q,h,k}(1).\label{eq:asymp}
\end{equation}
The exponential series, written with both signs before complexifying \(t\), converges normally on every bounded horizontal strip, together with all derivatives.
\end{proposition}
\begin{proof}
Gaussian inversion writes the prime-power sum twisted by \(\chi\) as
\[
 \frac1{2\pi i}\int_{\Re s=2}
 \frac{-L'(s,\chi)}{L(s,\chi)}
 \exp\{h(s-\tfrac12)^2+t(s-\tfrac12)\}\dd s.
\]
Move to \(\Re s=1/4\). The standard local partial-fraction formula and the unit-interval zero count \(O_q(\log(2+|u|))\) give a logarithmic bound for \(L'/L(1/4+iu,\chi)\) for large \(|u|\), since critical zeros have horizontal distance \(1/4\). Imprimitive Euler factors have zeros on \(\Re s=0\) and bounded logarithmic derivatives in the strip used here. Choose horizontal sides in \([N,N+1]\) and \([-N-1,-N]\) at distance \(\gg_q1/\log N\) from nearby ordinates. The logarithmic derivative on these sides is \(O_q(\log^2 N)\). These estimates follow from the Dirichlet \(L\)-function partial fractions in \cite{kouk}. The Gaussian makes the horizontal integrals tend to zero even after inserting \((s-1/2)^k\).
The new vertical line contributes \(O_{q,h,k}(\e^{-t/4})\). Each critical zero contributes its negative multiplicity times \(\exp(-h\gamma^2+i\gamma t)\), and the principal pole contributes \(\exp(h/4+t/2)\), before differentiation. Character orthogonality gives the first formula. The local zero count proves normal convergence.

The terms \(p^j\) with \(j\ge3\) have total absolute coefficient at most
\(\sum_p(\log p)p^{-3/2}/(1-p^{-1/2})<\infty\), so their translated Gaussian and each derivative tend to zero. For squares in class \(a\), each class \(b\) with \(b^2=a\) contributes
\[
 \sum_{p\equiv b\ (q)}\frac{\log p}{p}\,k_h^{(k)}(t-2\log p).
\]
The prime number theorem in progressions and partial summation show that the measure
\(\sum_{p\equiv b\ (q)}(\log p)p^{-1}\delta_{2\log p}\), shifted left by \(t\), gives mass tending to \((B-A)/(2d)\) on each fixed interval \([A,B]\). Indeed the endpoints for \(p\) are \(\exp((t+A)/2)\) and \(\exp((t+B)/2)\); substitute \(\vartheta(y;q,b)=y/d+o_q(y)\) in partial summation on that interval. The mass of every translated unit interval is bounded by the Chebyshev bound \(\vartheta(y)\ll y\). Gaussian decay therefore permits passage from interval tests to \(k_h^{(k)}(-u)\) by truncation and uniform approximation. Its integral is one for \(k=0\) and zero otherwise. Summing over square roots proves the second formula.
\end{proof}

\begin{lemma}[No constant direction]\label{le:rank}
For every real \(v\ne0\), some \(\gamma>0\) satisfies \(v\cdot b_\gamma\ne0\). More generally, if complex \(c_\chi\) satisfy \(\sum_\chi c_\chi m_\chi(\gamma)=0\) for every nonzero signed ordinate, all \(c_\chi\) vanish.
\end{lemma}
\begin{proof}
Use the entire order-one completion
\[
 \Xi_\chi(s)=[s(s-1)]^{\delta_\chi}
 (f_\chi/\pi)^{(s+a_\chi)/2}\Gamma((s+a_\chi)/2)L(s,\chi^*),
\]
where \(f_\chi\) is the conductor, \(a_\chi\in\{0,1\}\) the parity, and \(\delta_\chi=1\) for the conductor-one principal character only. Its zeros are exactly the nontrivial zeros, with multiplicity. The genus-one Hadamard product gives
\[
 \sum_\chi c_\chi\,\frac{\Xi_\chi'}{\Xi_\chi}(s)= A+\frac{M}{s-1/2},
 \qquad M=\sum_\chi c_\chi m_\chi(0).
\]
Indeed the grouped terms \(1/(s-\rho)+1/\rho\) converge normally off the zeros; all noncentral groups cancel. Put \(D(s)=\sum c_\chi L'(s,\chi^*)/L(s,\chi^*)\). Subtract the identity at \(s\) from that at \(s+2\). The conductor constants disappear and the logarithmic gamma derivatives differ by \(1/(s+a_\chi)\), including the factor \(1/2\) from the chain rule. The factors \(s(s-1)\) contribute rational terms. Hence \(D(s+2)-D(s)\) is rational. The Euler series gives \(D(\sigma)=O(2^{-\sigma})\) as \(\sigma\to+\infty\) for fixed coefficients. A rational function with this decay on the positive real axis vanishes identically. Iterating the period \(2\) gives \(D=0\) for real \(s>1\). Uniqueness of Dirichlet coefficients gives \(\sum c_\chi\chi(p)=0\) at every prime \(p\nmid q\). Each reduced class contains a prime, and the character table is invertible. For real \(v\), annihilation at positive ordinates implies annihilation at negative ordinates by conjugation. Apply the stronger assertion to \(c_\chi=\sum_a v_a\overline{\chi(a)}\) and invert the character table again.
\end{proof}
